\begin{tcolorbox}[
  colback=IEEEblue!5,
  colframe=IEEEblue!70!black,
  fonttitle=\bfseries\sffamily,
  toptitle=1mm,
  bottomtitle=1mm,
  boxrule=0.6pt,
  arc=0mm,
  left=3mm,right=3mm,top=1.5mm,bottom=1.5mm
]
\footnotesize
\setlength{\tabcolsep}{4pt}
\renewcommand{\arraystretch}{1.08}
\begin{tabular}{@{}>{\bfseries\RaggedRight\arraybackslash}p{0.20\linewidth}
                    >{\RaggedRight\arraybackslash}p{0.31\linewidth}
                    >{\RaggedRight\arraybackslash}p{0.41\linewidth}@{}}
  \toprule
  \textbf{Check} & \textbf{Correct Method} & \textbf{Common Mistake} \\
  \midrule
  Entropy & Compute $-\sum p_i\log_2p_i$ from measured frequencies & Reporting entropy as the actual Huffman length \\
  Huffman Length & Calculate $\sum f_i\ell_i$ from all codeword lengths & Assuming entropy is directly achievable per digit \\
  Compression Ratio & Use one clearly stated storage baseline & Mixing four-bit, ASCII, and packed-decimal baselines \\
  Storage Cost & Include codebook, metadata, and padding & Reporting only ideal payload bits \\
  Verification & Decode and compare with the original sequence & Treating encoding alone as proof of losslessness \\
  \bottomrule
\end{tabular}
\end{tcolorbox}